\section{Sentiment im Embedding-Raum}

Dieses Kapitel untersucht die erste und die vierte der fünf Methoden. Die Residual-Stream-Analyse betrachtet die geometrische Struktur der \glspl{token}-Repräsentationen, und das lineare Probing prüft, ob sich positive und negative Wörter linear trennen lassen. Beide Analysen setzen am Eingang des Modells an, also an den \glspl{activation} der Embedding-Schicht vor der ersten \gls{transformer}-Schicht. Grundlage sind die 2.254 Ein-Token-Wörter des Hu-\&-Liu-Lexikons.

\subsection{Geometrie der Sentiment-Embeddings}

\subsubsection{Hauptkomponentenanalyse}

Um die hochdimensionalen Embeddings sichtbar zu machen, wird eine Hauptkomponentenanalyse durchgeführt. Dieses Verfahren reduziert die 1.024 Dimensionen auf wenige Achsen, die möglichst viel der Streuung der Daten erfassen. Abbildung~\ref{fig:pca} zeigt die positiven und negativen Wörter entlang der ersten beiden Hauptkomponenten sowie ihre Verteilung entlang der ersten Achse.

\begin{figure}[htbp]
  \centering
  \includegraphics[width=0.49\textwidth]{hu_liu_opinion_embeddings.png}
  \hfill
  \includegraphics[width=0.49\textwidth]{hu_liu_opinion_histogram.png}
  \caption{Positive und negative Hu-\&-Liu-Wörter im Raum der ersten beiden Hauptkomponenten (links) und ihre Verteilung entlang der ersten Hauptkomponente (rechts).}
  \label{fig:pca}
\end{figure}

Die ersten beiden Hauptkomponenten trennen die Sentiment-Klassen nicht vollständig. Es zeigen sich jedoch räumliche Tendenzen, und positive sowie negative Wörter verteilen sich entlang der ersten Achse unterschiedlich. Die Überlappung deutet darauf hin, dass Sentiment nicht durch zwei reduzierte Dimensionen vollständig beschrieben wird, sondern über viele Dimensionen verteilt ist. Abbildung~\ref{fig:fieldpanels} zeigt die gleiche Betrachtung getrennt nach thematischen Feldern.

\begin{figure}[htbp]
  \centering
  \includegraphics[width=0.9\textwidth]{hu_liu_field_embedding_panels.png}
  \caption{Embeddings positiver und negativer Wörter getrennt nach thematischen Feldern.}
  \label{fig:fieldpanels}
\end{figure}

\subsubsection{Normen der Embeddings}

Die L2-Norm beschreibt die Länge eines Vektors. Abbildung~\ref{fig:norms} vergleicht die Normen positiver und negativer Wörter innerhalb einzelner Felder. Die absolute Spannweite ist mit Werten zwischen etwa 0,78 und 0,82 gering, innerhalb dieser Spannweite zeigt sich jedoch ein konsistentes Muster: In nahezu allen zehn Feldern liegt die Norm der negativen Wörter etwas über der Norm der positiven Wörter, etwa bei „sour“ gegenüber „delicious“, „dangerous“ gegenüber „safe“ oder „failed“ gegenüber „acclaimed“. Die Vektorlänge allein eignet sich wegen der geringen absoluten Unterschiede weiterhin nicht als praktikabler Indikator für Sentiment, der systematische Versatz deutet aber darauf hin, dass negative Wörter im Embedding-Raum von Pythia-410M tendenziell etwas größere Normen besitzen als positive. Aussagekräftiger für die Trennung von Sentiment sind die Richtungen der Vektoren und ihre Nachbarschaften.

\begin{figure}[htbp]
  \centering
  \includegraphics[width=0.85\textwidth]{sentiment_embedding_norms_all_fields.png}
  \caption{L2-Normen der Sentiment-Embeddings getrennt nach Feldern.}
  \label{fig:norms}
\end{figure}

\subsubsection{Nachbarn, Arithmetik und Sentiment-Richtung}

Eine Analyse der nächsten Nachbarn über die Cosine-Ähnlichkeit zeigt, dass die einem Zielwort am nächsten liegenden Token meist dieselbe Sentiment-Polarität besitzen. Positive Wörter sind überwiegend von positiven Nachbarn umgeben, negative Wörter von negativen. Dieses Muster gilt sowohl für das gesamte Vokabular als auch für die auf das Lexikon beschränkte Suche.

Auch einfache Rechnungen mit Embedding-Vektoren bestätigen diese Struktur. Wird etwa die Kombination aus mehreren positiven Wörtern gebildet und ein negatives Wort abgezogen, liegen die nächsten Nachbarn des Ergebnisses überwiegend im positiven Bereich. Umgekehrt führen Kombinationen mehrerer negativer Wörter zu überwiegend negativen Nachbarn. Sentiment verhält sich somit als eine kontinuierliche und weitgehend linear strukturierte Eigenschaft des Embedding-Raums.

Besonders deutlich wird dies durch eine Sentiment-Richtung, die aus der Differenz der Embeddings von „good“ und „bad“ gebildet wird. Werden die Sentiment-Wörter auf diese Richtung projiziert, ergibt sich eine klare Trennung. Positive Wörter besitzen überwiegend positive Projektionswerte, negative Wörter überwiegend negative. Stärker konnotierte Wörter liegen tendenziell weiter von der Mitte entfernt. Abbildung~\ref{fig:direction} stellt diese Projektion dar. Die Richtung „good“ minus „bad“ kann damit als eine Art Sentiment-Achse verstanden werden.

\begin{figure}[htbp]
  \centering
  \includegraphics[width=0.85\textwidth]{sentiment_direction_projection_all_fields.png}
  \caption{Projektion der Sentiment-Wörter auf die Richtung „good“ minus „bad“.}
  \label{fig:direction}
\end{figure}

\subsection{Lineares Probing der Embeddings}

Um die geometrischen Beobachtungen quantitativ zu prüfen, wird ein lineares Probing durchgeführt. Dabei wird eine logistische Regression auf den 1.024-dimensionalen Embeddings der Hu-\&-Liu-Wörter trainiert. Die logistische Regression ist ein linearer Klassifikator. Eine hohe Genauigkeit bedeutet, dass positive und negative Wörter bereits durch eine einfache lineare Entscheidungsgrenze getrennt werden können. Die Daten werden geschichtet und im Verhältnis von 80 zu 20 in Trainings- und Testdaten aufgeteilt. Tabelle~\ref{tab:split} zeigt die Aufteilung.

\begin{table}[htbp]
  \centering
  \caption{Geschichtete Aufteilung der Hu-\&-Liu-Wörter für das lineare Probing.}
  \label{tab:split}
  \begin{tabular}{lrrr}
    \toprule
    Teilmenge & positiv & negativ & gesamt \\
    \midrule
    Training & 692 & 1.111 & 1.803 \\
    Test & 173 & 278 & 451 \\
    \midrule
    gesamt & 865 & 1.389 & 2.254 \\
    \bottomrule
  \end{tabular}
\end{table}

Auf den Testdaten erreicht die lineare Probe eine Klassifikationsgenauigkeit von 97 Prozent bei gleichzeitig hohen Präzisions-, Recall- und F1-Werten für beide Klassen. Dieses Ergebnis zeigt, dass die Sentiment-Information im Embedding-Raum weitgehend linear separierbar ist, obwohl die ersten beiden Hauptkomponenten nur einen kleinen Teil der Streuung erklären. Die Information ist demnach über viele Dimensionen verteilt. Abbildung~\ref{fig:lraxis} und Abbildung~\ref{fig:lrfields} zeigen die durch die Probe gelernte Sentiment-Achse sowie die vorhergesagten Wahrscheinlichkeiten für einzelne Feldwörter.

\begin{figure}[htbp]
  \centering
  \includegraphics[width=0.8\textwidth]{logistic_regression_sentiment_axis.png}
  \caption{Durch die lineare Probe gelernte Sentiment-Achse.}
  \label{fig:lraxis}
\end{figure}

Abbildung~\ref{fig:lraxis} markiert die von der Probe falsch klassifizierten Wörter zusätzlich farblich. Die Fehlklassifikationen betreffen erkennbar Wörter, deren Sentiment stark vom Kontext oder von einer untergeordneten, fachsprachlichen Lesart abhängt: „capricious“ und „tout“ werden mit Wahrscheinlichkeiten von rund 0,88 beziehungsweise 0,81 als positiv vorhergesagt, obwohl beide Wörter im Lexikon als negativ gelten, „unsupported“ erhält ebenfalls einen hohen positiven Wert von etwa 0,82. In die andere Richtung wird „tempt“ mit nur rund 0,07 stark negativ eingeschätzt. Das Wort „morality“ liegt mit einer Wahrscheinlichkeit nahe 0,50 fast exakt auf der Entscheidungsgrenze, was zu der inhaltlichen Erwartung passt, dass moralische Bewertungen ohne weiteren Kontext kein eindeutiges Sentiment tragen. Diese Einzelfälle zeigen, dass die verbleibenden rund drei Prozent Fehlklassifikationen nicht zufällig verteilt sind, sondern gezielt bei mehrdeutigen oder kontextabhängigen Wörtern auftreten.

\begin{figure}[htbp]
  \centering
  \includegraphics[width=0.85\textwidth]{logistic_regression_field_word_probabilities.png}
  \caption{Von der linearen Probe vorhergesagte Sentiment-Wahrscheinlichkeiten für ausgewählte Feldwörter.}
  \label{fig:lrfields}
\end{figure}

Zusammengenommen liefern die geometrischen Analysen und die lineare Probe übereinstimmende Hinweise. Der Embedding-Raum von Pythia-410M enthält bereits vor der \gls{transformer}-Verarbeitung eine strukturierte Repräsentation von Sentiment. Diese zeigt sich in sentimenthomogenen Nachbarschaften, in sinnvollen Ergebnissen der Embedding-Arithmetik, in einer erkennbaren Sentiment-Achse und in einer hohen linearen Trennbarkeit. Damit ist die erste Forschungsfrage bejaht. Offen bleibt, ab welcher Schicht das Modell diese Information für seine Vorhersage tatsächlich nutzt.
